\section{集成学习概述}
\label{sec:ensemble-overview}

\subsection{从单个模型到组合预测}
\label{sec:ensemble-from-single-to-combination}

第\ref{sec:regression-trees}节和第\ref{sec:classification-trees}节已经说明，树的分裂由有限训练样本决定。某些候选分裂的得分很接近时，替换少量样本就可能改变根节点，后续分支也随之变化。单棵树因而可能具有较大的训练波动。另一方面，一个受限的线性模型即使训练稳定，也可能无法表示风险只在局部条件组合下升高的现象。前一种困难主要来自估计不稳定，后一种困难主要来自表示受限；反复调节同一个模型的复杂度，不一定能同时消除二者。

集成学习在原有预测任务之外增加两个环节：先产生多个\term{成员模型}（member model），再组合这些成员的输出。强调训练方法或模型族时，也常把产生成员的算法称为\term{基学习器}（base learner）。成员可以是同一算法在不同样本或随机种子下的结果，也可以来自不同模型族；最终组合得到的对象称为\term{集成模型}（ensemble model）。这种做法没有改变响应或类别标签的含义，只改变了由训练数据形成预测器的路径。

组合方法围绕三个问题逐渐分化。Bagging从训练集重采样多个成员并固定聚合规则，使易受数据扰动影响的预测更稳定；Boosting让下一轮成员依赖当前组合，把新增容量分配给尚未处理好的训练信号；Stacking则用样本外成员预测训练另一层模型，使组合规则也可以学习。

图\ref{fig:ensemble-three-mechanisms}固定同一训练任务，比较三类方法改变的环节。Bagging中的成员可以彼此独立训练；Boosting的第$t$个成员依赖前$t-1$轮形成的当前预测；Stacking的基成员可以独立训练，但组合层必须看到不含当前样本训练信息的折外预测。这些差异比“用了多少个模型”更能说明训练成本与泄漏风险。

\begin{figure*}[htbp]
  \centering
  % Mechanism comparison for Bagging, Boosting, and Stacking.
% Schematic only; no empirical quantities are encoded.
% Canvas: 164 mm x 66 mm. Dependencies: project palette and TikZ.
\begin{tikzpicture}[
  x=1mm,y=1mm,
  font=\figurefont\footnotesize,
  text=FigureInk,
  source/.style={
    rectangle,
    model input,
    minimum width=29mm,
    minimum height=11mm,
    text width=25mm,
    align=center,
    inner sep=1.5mm
  },
  process/.style={
    rectangle,
    model hidden,
    minimum width=32mm,
    minimum height=11mm,
    text width=28mm,
    align=center,
    inner sep=1.5mm
  },
  combine/.style={
    rectangle,
    model training,
    minimum width=30mm,
    minimum height=11mm,
    text width=26mm,
    align=center,
    inner sep=1.5mm
  },
  output/.style={
    rectangle,
    model output,
    minimum width=18mm,
    minimum height=11mm,
    align=center,
    inner sep=1.5mm
  },
  flow/.style={
    model flow
  },
  dependency/.style={
    model feedback
  }
]
  \path[use as bounding box] (0,0) rectangle (164,66);

  \node[anchor=west,font=\figurefont\bfseries\small,text=FigureInk] at (1,55) {Bagging};
  \node[source] (bag-source) at (40,55) {重采样或随机化};
  \node[process] (bag-members) at (82,55) {并行成员\\$h_1,\ldots,h_M$};
  \node[combine] (bag-combine) at (123,55) {固定平均或投票};
  \node[output] (bag-output) at (154,55) {$H$};
  \draw[flow] (bag-source) -- (bag-members);
  \draw[flow] (bag-members) -- (bag-combine);
  \draw[flow] (bag-combine) -- (bag-output);

  \node[anchor=west,font=\figurefont\bfseries\small,text=FigureInk] at (1,33) {Boosting};
  \node[source] (boost-source) at (40,33) {当前组合$F_{t-1}$};
  \node[process] (boost-members) at (82,33) {序贯训练成员$h_t$};
  \node[combine] (boost-combine) at (123,33) {加法更新$F_t$};
  \node[output] (boost-output) at (154,33) {$F_T$};
  \draw[dependency] (boost-source) -- node[above=0.8mm,text=FigureInk] {修正信号} (boost-members);
  \draw[flow] (boost-members) -- (boost-combine);
  \draw[flow] (boost-combine) -- (boost-output);
  \draw[dependency] (boost-combine.south) -- ++(0,-5)
    -| node[pos=0.35,above=0.8mm,text=FigureInk] {下一轮} (boost-source.south);

  \node[anchor=west,font=\figurefont\bfseries\small,text=FigureInk] at (1,11) {Stacking};
  \node[source] (stack-source) at (40,11) {折内训练基成员};
  \node[process] (stack-members) at (82,11) {折外预测矩阵$\bm{Z}$};
  \node[combine] (stack-combine) at (123,11) {训练元学习器$g$};
  \node[output] (stack-output) at (154,11) {$H$};
  \draw[flow] (stack-source) -- (stack-members);
  \draw[flow] (stack-members) -- (stack-combine);
  \draw[flow] (stack-combine) -- (stack-output);

  \draw[FigureGrid,line width=0.6pt] (0,44)--(164,44);
  \draw[FigureGrid,line width=0.6pt] (0,22)--(164,22);
\end{tikzpicture}

  \caption{Bagging、Boosting与Stacking的成员生成及组合关系。实线表示训练或预测的主路径，虚线表示由前一轮组合产生的修正信号。三类方法都输出组合预测，但成员间依赖与组合规则不同；图中流程为机制示意，不表示三者必须使用相同成员模型。}
  \label{fig:ensemble-three-mechanisms}
\end{figure*}

\subsection{成员、组合规则与输出}
\label{sec:ensemble-members-combination-output}

设训练集为
\[
  D_n=\{(\bm{x}_i,y_i)\}_{i=1}^{n},
\]
学习过程产生$M$个成员模型$h_1,\ldots,h_M$。用$\mathcal{C}$表示组合规则，集成预测统一写成
\begin{equation}
  H(\bm{x})
  =\mathcal{C}\bigl(h_1(\bm{x}),\ldots,h_M(\bm{x})\bigr).
  \label{eq:ensemble-combination}
\end{equation}
式\eqref{eq:ensemble-combination}只规定组合接口，还没有规定输出的语义。回归成员若都预测同一个数值目标，可以采用等权平均
\[
  H(\bm{x})=\frac{1}{M}\sum_{m=1}^{M}h_m(\bm{x}),
\]
也可以用非负权重$w_m$作加权平均，其中$\sum_m w_m=1$。分类成员则可能输出硬标签、类别得分或类别概率。硬标签适合投票，得分和概率可以逐类平均后再决策，但三种量不能在没有转换的情况下混合。

输出类型之所以重要，是因为相同的代数操作可能对应不同结论。两个分类器的得分尺度不同，直接平均会让数值范围较大的成员取得更大影响；两个成员都输出概率，也不意味着平均后已经校准。若元学习器把成员得分映射为概率，还要单独检查这个映射在新样本上的可靠性。第\ref{sec:classification-evaluation}节给出的硬标签、排序得分和类别概率评价接口同样适用于集成模型。

组合权重也不是成员的固有属性。等权平均预先固定权重，AdaBoost在序贯训练中同时确定成员贡献，线性Stacking则根据折外预测估计权重。权重还可以依赖输入，此时$w_m=w_m(\bm{x})$，不同样本会调用不同成员或采用不同组合比例。成员数量、权重约束、组合层复杂度以及是否允许输入相关组合，都属于训练或验证流程需要确定的模型选择。

\subsection{成员质量、错误依赖与多样性}
\label{sec:ensemble-quality-dependence-diversity}

平均为什么可能比单个成员稳定，可以先看一个受控情形。固定输入$\bm{x}$，把训练集抽取、重采样和训练随机性都视为随机来源。若$M$个成员预测具有相同方差$\sigma^2(\bm{x})$，任意两个不同成员的预测相关系数均为$\rho(\bm{x})$，则等权平均的方差为
\begin{equation}
  \operatorname{Var}\!\left[H(\bm{x})\right]
  =
  \sigma^2(\bm{x})
  \left[
    \rho(\bm{x})
    +
    \frac{1-\rho(\bm{x})}{M}
  \right].
  \label{eq:ensemble-correlated-variance}
\end{equation}
其中随$M$衰减的是成员不共享的波动；$\rho(\bm{x})\sigma^2(\bm{x})$所代表的共同变化不会因继续增加同类成员而消失。式\eqref{eq:ensemble-correlated-variance}依赖等方差、等相关的简化设定，用来解释相关性的作用，不是任意集成的精确风险公式。

图\ref{fig:ensemble-error-dependence}把这一差别写成成员误差的两部分。成员特有波动$u_m$在平均中可能相互抵消；若每个成员都带有共同误差$c$，增加同类成员只能继续削弱前一部分，不能消除$c$。因此，成员数增长与成员错误依赖下降是两个不同的条件。

\begin{figure*}[htbp]
  \centering
  % Averaging suppresses member-specific fluctuations but preserves shared error.
% Schematic only; no empirical quantities are encoded.
% Canvas: 169 mm x 58 mm. Dependencies: project palette and TikZ.
\begin{tikzpicture}[
  x=1mm,y=1mm,
  font=\figurefont\footnotesize,
  text=FigureInk,
  member/.style={
    circle,
    model hidden,
    minimum size=8mm,
    inner sep=0.5mm
  },
  shared/.style={
    rectangle,
    model training,
    minimum width=50mm,
    minimum height=8mm,
    align=center,
    inner sep=1mm
  },
  result/.style={
    rectangle,
    model output,
    minimum width=54mm,
    minimum height=10mm,
    text width=50mm,
    align=center,
    inner sep=1mm
  },
  flow/.style={
    model flow
  }
]
  \path[use as bounding box] (0,0) rectangle (169,58);

  \node[anchor=west,font=\figurefont\small] at (1,54)
    {(a) 只有成员特有波动};
  \foreach \x/\name/\lab in {
    8/u1/$u_1$,21/u2/$u_2$,34/u3/$u_3$,47/u4/$u_4$,60/u5/$u_M$} {
    \node[member] (\name) at (\x,38) {\lab};
  }
  \node[result] (independent-average) at (34,12)
    {平均后只剩\\$\displaystyle M^{-1}\sum_m u_m$};
  \foreach \name in {u1,u2,u3,u4,u5} {
    \draw[FigureInk,line width=.8pt] (\name.south)--++(0,-6);
  }
  \draw[FigureInk,line width=.8pt] (8,28)--(60,28);
  \draw[flow] (34,28)--(independent-average.north);

  \draw[FigureGrid,line width=0.6pt] (82,3)--(82,55);

  \node[anchor=west,font=\figurefont\small] at (87,54)
    {(b) 成员还共享同一偏移};
  \node[shared] (common) at (126,45)
    {共同误差$c$进入每个成员};
  \node[member] (cu1) at (92,30) {$u_1$};
  \node[member] (cu2) at (107,30) {$u_2$};
  \node[member] (cu4) at (145,30) {$u_{M-1}$};
  \node[member] (cu5) at (160,30) {$u_M$};
  \node[result] (correlated-average) at (126,9)
    {平均误差\\$\displaystyle c+M^{-1}\sum_m u_m$};
  \foreach \name in {cu1,cu2,cu4,cu5} {
    \draw[FigureInk,line width=.8pt] (\name.south)--++(0,-6);
  }
  \draw[FigureInk,line width=.8pt] (92,20)--(160,20);
  \draw[flow] (126,20)--(correlated-average.north);
  \draw[model feedback] (common.east)--(166,45)
    |-(correlated-average.east);
\end{tikzpicture}

  \caption{成员误差依赖与平均的作用。左图表示只有成员特有波动时，平均可以随成员数增加而削弱这些波动；右图表示所有成员还共享误差$c$时，共同项仍保留在平均预测中。图中成员数量与符号只用于说明机制，不表示实测误差。}
  \label{fig:ensemble-error-dependence}
\end{figure*}

平方误差还给出另一个逐点恒等式。对非负权重$w_m$且$\sum_mw_m=1$，令$H(\bm{x})=\sum_mw_mh_m(\bm{x})$，则
\begin{equation}
  \bigl(y-H(\bm{x})\bigr)^2
  =
  \sum_{m=1}^{M}w_m\bigl(y-h_m(\bm{x})\bigr)^2
  -
  \sum_{m=1}^{M}w_m\bigl(h_m(\bm{x})-H(\bm{x})\bigr)^2.
  \label{eq:ensemble-ambiguity-decomposition}
\end{equation}
右侧第一项是成员平方误差的加权平均，第二项衡量成员围绕组合预测的分散程度，常称为歧义项\scite{krogh1995ensembles}在成员平均误差固定时，更大的有效分散会降低组合误差；但若为了制造分散而加入很差的成员，第一项也会增大。多样性只能与成员质量一起解释。

分类投票没有完全相同的平方误差分解，却面对同一问题：关键不只是每个成员错多少次，还包括它们是否在同一批样本上出错。对一对分类器，可以记录二者同时正确、只有一个正确和同时错误的样本数$n_{1,1},n_{1,0},n_{0,1},n_{0,0}$。分歧率$(n_{1,0}+n_{0,1})/n$描述二者给出不同正确性结果的频率，双错率$n_{0,0}/n$直接关注共同错误；Q统计量
\[
  Q=
  \frac{n_{1,1}n_{0,0}-n_{1,0}n_{0,1}}
       {n_{1,1}n_{0,0}+n_{1,0}n_{0,1}}
\]
在分母非零时描述两种正确性结果的关联。不同指标关注的对象不同，也不存在一个脱离成员准确率、对所有任务都能单调预测集成风险的多样性分数\scite{kuncheva2003diversity}

成员差异可以从多个位置产生。改变训练样本会改变成员看到的经验分布；改变输入特征会改变可用信息；改变模型结构或超参数会改变表示与正则；改变初始化、样本顺序或分裂候选会改变优化路径。Bagging主要扰动训练样本，随机森林和Extra Trees进一步扰动特征或阈值，Stacking则允许直接组合不同模型族。这些机制只是多样性的来源，最终是否形成互补仍要用样本外预测检查。

由此可以用三个维度比较后续方法：成员怎样产生，成员之间有什么依赖，组合规则怎样确定。Bagging采用扰动后的成员和固定聚合，Boosting采用序贯依赖成员和加法组合，Stacking采用样本外成员预测和学习式组合。三类方法都可能改善预测，也都可能因共同失配、数据泄漏或过度选择而失效。
